%% =========================================================================
%%  Заключение
%% =========================================================================

\chapter*{Заключение}
\addcontentsline{toc}{chapter}{Заключение}

\section*{Основные результаты работы}
\addcontentsline{toc}{section}{Основные результаты работы}

В работе развита единая точка зрения на квазиньютоновские обновления
с памятью прошлых секущих условий и последовательно применена
к задачам нелинейных систем и гладкой оптимизации. Ключевые
результаты по главам перечислены ниже.

\paragraph{Глава 1 (приближение якобиана).}
Выведена общая формула семейства мультисекантных обновлений
\[
  B_{k+1} = \widehat{B}_k
    + (Y_m - \widehat{B}_k S_m)\,(V_m^\top S_m)^{-1}\,V_m^\top,
\]
охарактеризующая все ранг-$m$ обновления якобиана, которые точно
удовлетворяют заданному набору секущих условий. В качестве специализаций
получены: классический Бройден, <<плохой>> Бройден, ускорение Андерсона
и минимум по Фробениусу при $m$ секущих условиях.
Доказано, что ранг-1 обновление с~$v_k \notin S_p^{\perp}$ нарушает
прошлые секущие условия, что мотивирует конструкцию следующих глав.

\paragraph{Глава 2 (мультисекущий Бройден).}
Рассматривается мультисекущее обновление Бройдена~\cite{schnabel1983}
\[
  B_{k+1} = B_k + (Y_p - B_k S_p)\,G_p^{-1}\,S_p^\top,
  \qquad G_p = S_p^\top S_p,
\]
$S_p$, $Y_p$~--- окно текущего и $p$ прошлых секущих пар. Ранг-1
специализация под индуктивным инвариантом $B_k s_j = y_j$ ($j<k$)
совпадает с Projected Broyden Гэя--Шнабеля~\cite{gay_schnabel1978} и
используется в~\S\ref{sec:highdim} для limited-memory варианта L-PB
через формулу Шермана--Моррисона.

Основной результат главы~--- \textbf{теорема об ограниченной деградации
с расширенным проектором} (\ref{thm:bd}): точное \emph{равенство}
\[
  \|E_{k+1}\|_F^2 = \|E_k\|_F^2 - \|E_k \Pi_k\|_F^2 + \Delta_k
\]
с проектором $\Pi_k$ ранга $p{+}1$ на $\operatorname{col}(S_p)$
(у классического Бройдена~--- ранга~1, на $\operatorname{span}(s_k)$)
и явной геометрической константой
$C_{\mathrm{PB}} = L(1+C_s/2)\sqrt{p+1}\,\kappa_p$, где
$\kappa_p = \kappa(S_p)$~--- обусловленность окна шагов.
У Гэя--Шнабеля для ранг-1 специализации
оценка bounded deterioration установлена в виде неравенства сверху
(Th.~4.2); здесь равенство с явной геометрической константой и
переход к мультисекущей форме сделаны впервые.
Доказано неравенство
$\|E_k\Pi_k\|_F^2 \ge \|E_k s_k\|^2/\|s_k\|^2$: отрицательный член
снимает с матрицы ошибки \emph{не меньше} информации, чем у
классического Бройдена, и сохраняются $p{+}1$ секущих условий
$B_{k+1} s_j = y_j$, $j = k{-}p, \ldots, k$.

В качестве следствия установлена количественная
\textbf{$O(\rho_k)$-аппроксимация якобиана на подпространстве прошлых
шагов} $\operatorname{col}(S_p)$ (\ref{cor:subspace-alignment})~---
структурное преимущество над классическим Бройденом, у которого
аппроксимация контролируется лишь вдоль одномерного направления $s_k$.

На задаче Discrete~BVP ($n{=}100$) мультисекущий Бройден работает
устойчиво там, где классический Бройден неустойчив или сходится
медленнее. Limited-memory вариант L-PB при $m\ge p_{\max}+1$
устойчиво сходится на Broyden Banded при $n\in\{10^4, 10^5\}$ за~число
итераций, не~зависящее от~$n$; зазор с~плотной формой умеренный и~монотонно
сокращается с~ростом~$m$, что отражает типичную для~limited-memory методов
картину (ср.\ L-BFGS).

\paragraph{Глава 3 (SS-U и SS-SR1).}
Воспроизведён классический результат Шнабеля~\cite{schnabel1983}
о несовместности точного сохранения $p$~прошлых секущих условий
с~требованием симметрии (лемма~\ref{lem:incompat}): систему
одновременно нельзя удовлетворить в общем положении.
Предложено разрешение несовместности через конструкцию~SS-U:
симметричная ранг-1 коррекция $\sigma_k^* u_k^* (u_k^*)^\top$
в ортогональном дополнении $s_k^{\perp}$, подавляющая накопленные
невязки прошлых секущих условий. Установлены:
\begin{itemize}[leftmargin=*]
  \item точное выполнение текущего секущего условия и симметрия
        (\ref{thm:ss_props});
  \item монотонное убывание невязки в~$s_k^{\perp}$ по оптимальному
        направлению $u_k^*$;
  \item bounded deterioration в~композитной форме для $\mathrm{SS}\text{-}\mathcal{U}$
        под Lipschitz-непрерывным гессианом и~входной Q-линейной
        сходимостью (Лемма~\ref{lem:bd_composite});
  \item сверхлинейная сходимость по критерию Dennis--More
        (\ref{thm:ss_dm});
  \item эмпирически наблюдаемая устойчивая сходимость на семействе
        Curved~Valley (см.~\S\ref{sec:ss_numerics}).
\end{itemize}

\section*{Связь результатов}
\addcontentsline{toc}{section}{Связь результатов}

Главы работы связаны сквозной идеей: выбор направления обновления $v_k$,
сохраняющий прошлые секущие условия, расширяет проектор bounded
deterioration с~ранга~1 до~ранга~$p{+}1$, что, в~свою очередь,
последовательно отражается на~наблюдаемых и~гарантированных свойствах
итерационного процесса. В~нелинейных системах (глава~2) это даёт
усиленную оценку bounded deterioration с~явной геометрической константой
$C_{\mathrm{PB}} = L(1{+}C_s/2)\sqrt{p+1}\,\kappa_p$, зависящей
от~обусловленности окна~$\kappa_p$, а~не~от размерности~$n$;
в~оптимизации (глава~3) идея переносится через композитное обновление
SS-$\mathcal{U}$ с~bounded deterioration в~форме Леммы~\ref{lem:bd_composite}
и~условной Q-сверхлинейной сходимостью по~Денниса--Море. Расширение
проектора с~ранга~1 (классический Бройден) до~ранга~$p{+}1$
(мультисекущий Бройден) вместе с~сохранением $p{+}1$ секущих условий~---
содержательное уточнение классического анализа, дающее на~каждой
итерации строго не~меньше информации об~ошибке аппроксимации якобиана.

В~литературе закрыты следующие пробелы: количественная форма bounded
deterioration для общей мультисекущей формы Шнабеля~\cite{schnabel1983}
(у~Гэя--Шнабеля~\cite{gay_schnabel1978}~--- только неравенство сверху
для ранг-1 специализации); явная зависимость константы оценки от~геометрии
окна~$\kappa_p$, не~от~$n$~--- основание для применимости в~высокой
размерности; перенос идеи сохранения прошлых секущих на~симметричный
случай через коррекцию в~$s_k^\perp$ с~bounded deterioration в~композитной
форме (Лемма~\ref{lem:bd_composite}). Полученные количественные оценки
ошибки $\|B_k - J(x_k)\|$ входят константой в~оценки скорости
сходимости методов более высокого уровня, где аппроксимация якобиана
используется как технический блок~\cite{viqa2024}.

\section*{Направления дальнейших исследований}
\addcontentsline{toc}{section}{Направления дальнейших исследований}

\begin{itemize}[leftmargin=*]
  \item \textbf{Высокая размерность и limited memory.}
        L-PB-вариант (\S\ref{sec:highdim}), построенный на
        ранг-1 специализации с~Шерманом--Моррисоном, проверен на
        Broyden~Banded при $n\in\{10^4, 10^5\}$. Дальнейшее
        направление~--- сравнение с L-BFGS и preconditioned Anderson
        на расширенном тестовом наборе.
  \item \textbf{Обобщение SS-U на не гладкие задачи.}
        Построить негладкий аналог SS-U (proximal-SS-U) для
        композитных задач $\min f(x) + g(x)$ и исследовать его
        скорость на задачах с разреженной регуляризацией.
  \item \textbf{Практическая валидация на прикладных задачах.}
        Применить мультисекущий Бройден к задачам
        обучения (нелинейные системы, возникающие в
        физически-информированных нейронных сетях), а~также к
        равновесным моделям из экономики.
  \item \textbf{Глобализация и сравнение стратегий шага.}
        Формальный анализ глобализации мультисекущего Бройдена и
        SS-$\mathcal{U}$ (backtracking-Армихо, условия Вольфа,
        trust-region-модификации) и систематическое численное
        сравнение этих стратегий на широком тестовом наборе
        Mor\'e--Garbow--Hillstrom~--- естественное продолжение
        локальной теории, развитой в~гл.~\ref{ch:sp-broyden-theory}
        и~гл.~\ref{ch:symmetric}.
  \item \textbf{Корректная реализация и сравнение с современным Anderson.}
        Переделать реализацию Anderson-acceleration из~\cite{walker2011}
        с релаксацией и мультисекантной подгонкой в~$L^2$-норме и
        провести честное сравнение с мультисекущим Бройденом на всём
        тестовом наборе Broyden.
\end{itemize}
